\section{Cross-word counting and the original ordinary column slopes}
\label{sec:global-count}

\status{GVC1--GVC7 prove the original logarithmic \texttt{CountF}
package, including comparison of different words on the same datum.
GVC8--GVC11 prove the full \texttt{NumsL}/\texttt{WtsL} supplier,
retaining the derivative hypothesis required by its literal statement.
GVC12--GVC17 install the original ordinary column slopes and
\texttt{SlopesMF}. GVC18 shows that the literal
\texttt{FamP.sizes} field already supplies the eventual derivative
bound, so no additional derivative premise is needed on that family.
All results concern a supplied
original family; they do not construct its compatible selector or
second layer.}

Fix the original family $\mathcal F=\mathrm{Fm}$, with
$0<\epsilon<1/52$. At a member $x,i$ write $D=\mathcal F.D(x,i)$
and $\mathcal L=\mathcal F.L(x,i)$. The scalar and rank notation is
\[
 \begin{gathered}
 \ell=x^\epsilon,\quad b=e^\ell,\quad a=xb,\quad
 X'=\exp(x^{\epsilon/3}),\quad b'=a/X',\\
 w''=\frac{b'}{xn},\quad S=\mathcal L.\mathrm{Sc}=4k\Delta,\quad
 \eta=\mathcal L.\mathrm{etaSz},\\
 \mathcal A=\mathcal L.\mathrm{La}=\log(a\,\operatorname{rd}E).
 \end{gathered}\tag{GVC.1}
\]
Here $\Delta=D.s.\mathrm{Del}$, $f=[E:\mathbb Q]/2=2n$,
$N=k+m(t+1)$, and $n_T=\#D.\mathrm{host}.\mathrm{TT}$.
These are equalities in the original records. In particular the
parameter $\Delta$ is not chosen anew in this chapter.
A word is denoted $\omega$, to distinguish it from its
Harish--Chandra coordinates $u_{\omega,\tau,j}$.
All logarithms of arithmetic quantities below are of positive numbers.

The finite-family predicates use eventual thresholds before the member
$i$, the word $\omega$, the archimedean place $\tau$ and any
coefficient place. The auxiliary constant
$C_Y=\max\{5,C_{YL}+1\}$ below is unrelated to the fixed global
constant $C_0=3$. The exact GL4 maximum-slope coefficient remains
$C_5=4$; the free constants in \texttt{Conclusion5} and
\texttt{Conclusion6} will be denoted $C_-$ and $C_+$.

\subsection{The logarithmic Vandermonde comparison}

% BEGIN GVC1
\paragraph{GVC1: logarithmic perturbations and the harmonic bound.}
\label{claim:GVC1}
If $c>0$, $z,z'\ge c$ and $|z-z'|\le e$, then
\[
                  |\log z-\log z'|\le e/c.                 \tag{GVC.2}
\]
Indeed, if $z\ge z'$, integrate $1/s\le1/c$ over
$[z',z]$; interchange the two variables otherwise.
For an integer $q\ge1$,
\[
 H_q=\sum_{j=1}^{q}\frac1j
 \le1+\int_1^q\frac{ds}{s}=1+\log q.
\]
The integral comparison follows termwise from
$1/j\le\int_{j-1}^j ds/s$ for $j\ge2$.
Consequently, for $0\le r<k$,
\[
                  \sum_{s=r+1}^{k-1}\frac1{s-r}
                         \le1+\log k.                     \tag{GVC.3}
\]
An empty sum is zero, so this includes $r=k-1$ and $k=1$.
These prove the source's \texttt{abs\_log\_sub\_le\_div}
and \texttt{harmonic\_gap\_sum\_le} independently.
% END GVC1

% BEGIN GVC2
\paragraph{GVC2: separated coordinate intervals give the exact variation.}
\label{claim:GVC2}
Let $\Delta>0$ and let $u,u'$ have $k$ real coordinates. Suppose
that every pair $0<s<k$ satisfies
\[
 u_0-u_s,\ u'_0-u'_s\ge4k\Delta,\qquad
 |(u_0-u_s)-(u'_0-u'_s)|\le\Delta,
\]
and that every $1\le r<s<k$ satisfies
\[
 \begin{split}
 u_r-u_s,\ u'_r-u'_s&\ge3(s-r)\Delta,\\
 |(u_r-u_s)-(u'_r-u'_s)|&\le2\Delta.
 \end{split}
\]
Writing $V_k(u)=\prod_{r<s}(u_r-u_s)$, we have
\[
 \boxed{\quad
 |\log V_k(u)-\log V_k(u')|
 \le\frac14+\frac23 k(1+\log k).
 \quad}                                                     \tag{GVC.4}
\]
For $k=0$, both Vandermondes are the empty product $1$,
so the conclusion is immediate. For $k\ge1$ all their factors
are positive. Expand their logarithms and apply GVC1.
Each pair with $r=0$ costs at most $1/(4k)$, and their total is
at most $1/4$. For any fixed $r\ge1$ the cost is at most
$(2/3)\sum_{s>r}(s-r)^{-1}\le(2/3)(1+\log k)$.
There are at most $k$ such rows. The triangle inequality proves
(GVC.4), with exactly the source's coefficient $2/3$.
% END GVC2

% BEGIN GVC3
\paragraph{GVC3: check both endpoints and every interior interval on the original data.}
\label{claim:GVC3}
For the moment suppose $k\ge44$, $\Delta\ge1$, and at a given
place $\tau$ the actual coordinates satisfy
\[
 \begin{gathered}
 v_j-v_{j+1}=4\Delta,\qquad v_j\ge0,\\
 u_{\omega,j}=v_{j-1}-\tfrac12-\xi_{\omega,j}\Delta,
 \qquad \xi_{\omega,j}\in\{0,1\},\quad1\le j\le k-2,\\
 u_{\omega,0}-v_0\ge4k\Delta.
 \end{gathered}                                               \tag{GVC.5}
\]
At a non-top place assume
$u_{\omega,k-1}=v_{k-2}-1/2$; at a top place assume
$u_{\omega,k-1}\le-4k\Delta$.
Finally, require that both endpoints are independent of $\omega$.
These are exactly the parts of \texttt{DataA} used by the source
variation argument. They imply the hypotheses of GVC2 for any two
words on this same $D,\tau$.

Here is the complete pair check. The chain identity gives
$v_j=v_0-4j\Delta$. Each interior coordinate lies in the same
interval
\[
 [v_0-4(j-1)\Delta-\tfrac12-\Delta,\
   v_0-4(j-1)\Delta-\tfrac12]
\]
of length $\Delta$, for every word. Every coordinate after the
first is at most $v_0$, including the last in either place type.
The fixed first endpoint therefore gives the first-row lower bound
$4k\Delta$, while the difference of the two corresponding gaps
has absolute value at most $\Delta$.

For two interior coordinates at distance $d=s-r\ge1$, each gap
is at least $(4d-1)\Delta\ge3d\Delta$.
The two gaps differ by at most $2\Delta$.
If $s=k-1$ is non-top, substitute
$v_{k-2}-1/2$ for the last coordinate and obtain the identical
$(4d-1)\Delta$ estimate; its variation is at most $\Delta$.
If $s=k-1$ is top, nonnegativity of $v_{r-1}$ gives
$u_{\omega,r}\ge-\Delta-1/2$, and hence the gap is at least
$(4k-1)\Delta-1/2$. This dominates
$3(k-1-r)\Delta$: their difference is
$(k+2+3r)\Delta-1/2>0$.
The last endpoint is fixed, so here too the variation is at most
$\Delta$. This exhausts the pairs.

For the literal family, DSZ4 proves the cloud formulas and the
non-top endpoint formulas in (GVC.5) on the original
$\Delta=D.s.\mathrm{Del}$. The identities
$D.\mathrm{hcW\_first}$ and $D.\mathrm{hcW\_last}$ identify both
endpoints with the endpoints of the one original $D.\mathrm{hc}$,
so they are independent of the word, including at top places.
At a top place these endpoints are
\[
 u_{\omega,0}=a_T+(k-1)/2,\qquad
 u_{\omega,k-1}=-b_T-(k-1)/2.
\]
The original parameters give $kS\le w''\le b'\le b$ and
$a_T\ge a$, $b_T\ge b'$. Thus the last endpoint is at most
$-b'\le-S=-4k\Delta$. Also $v_0\le S$ and
$a=xb\ge2S$ once $x\ge2$, so $u_{\omega,0}-v_0\ge S$.
This verifies the top inequalities rather than assuming an
unproved full \texttt{RowsA} package.

DSZ4 also gives $\Delta\ge1$ and $m\ge x^{1/4}$,
with $k\ge m$. Enlarge one threshold so that $m\ge44$.
Every check above then holds for all $i,\omega,\tau$.
In particular (GVC.4) proves the literal
\texttt{vandermonde\_variation\_at\_place} whenever
\texttt{DataA} is supplied, and proves its needed actual-family
conclusion directly without claiming every stronger
\texttt{DataA}/\texttt{FamilyPartFive} clause.
% END GVC3

% BEGIN GVC4
\paragraph{GVC4: cancellation of the word-independent mass factors.}
\label{claim:GVC4}
The exact source definitions give
\[
 d_*^{\,\mathrm{form}}(\omega)
 =K_D\prod_{\tau}V_k(u_{\omega,\tau}),\qquad
 K_D=\frac{\mathrm{massI}(D,k)\,
          \mathrm{idxPi}(D,g_0,g_1)}
          {\prod_{\tau}\prod_{j=0}^{k-1}j!}.                 \tag{GVC.6}
\]
Both $K_D$ and $D_*^{\,\mathrm{form}}=\mathrm{DstarD}(D)$
belong to this same $D$ and are independent of $\omega$.
For the original
\[
 X_\omega=\frac1{fk}
          \log\frac{D_*^{\,\mathrm{form}}}
                    {d_*^{\,\mathrm{form}}(\omega)}
            =\mathrm{logXuD}(D,\omega)
\]
their contributions cancel in a difference. Since the original
place set has cardinality $f$, GVC3 gives
\[
 \boxed{\quad
 |X_\omega-X_{\omega'}|
 \le \frac1{4k}+\frac23(1+\log k)
 \le 1+\log k+\frac2k .
 \quad}                                                      \tag{GVC.7}
\]
Indeed expansion of the positive Vandermonde products bounds the
absolute value of their log difference by
$f[1/4+(2/3)k(1+\log k)]$; divide by $fk$.

This also proves the full generic
\texttt{VandermondeVariationExplicitBound} under \texttt{DataA}.
Its first requested bound is
$k(1+\log k)+2$, which is larger than (GVC.4).
For completeness, the formal generic expression allows a zero
common factor or a zero numerator. With the formal convention
$\log0=0$, and real logarithm $\log|z|$ for $z\ne0$, either zero
case makes both log-ratios zero. In the other cases
$\log|ab|=\log|a|+\log|b|$ gives the same cancellation.
No nonzero hypothesis is omitted in this generic statement.
For the actual family all these factors are positive by
QW1--QW2, the mass identity and the original positive degree fields,
so only the ordinary positive-logarithm calculation is used.

Finally the stored original equalities
$\mathcal F.\mathrm{ex}.\mathrm{duStar\_def}$ and
$\mathcal F.\mathrm{ex}.\mathrm{Dstar\_def}$ give
\[
 d_*^{\,\mathrm{actual}}(\omega)=d_*^{\,\mathrm{form}}(\omega),
 \qquad D_*^{\,\mathrm{actual}}=D_*^{\,\mathrm{form}}.
\]
Consequently $X_\omega=\log((D.\mathrm{mem}\,\omega).A.\mathrm{Xu})$.
This is the complete \texttt{jnt\_logXu} dictionary: the positive
real power in \texttt{Xu} has exponent $1/(fk)$.
No new mass, row degree, coefficient field or column is substituted.
% END GVC4

% BEGIN GVC5
\paragraph{GVC5: the same prime center gives cross-word prime variation.}
\label{claim:GVC5}
The original field $D.\mathrm{GamP\_eq}$ identifies all
$\Gamma_\omega=(D.\mathrm{mem}\,\omega).A.\mathrm{GamP}$
with the same positive $\Gamma$. Thus
\[
 \log T_\omega^*
       =X_\omega-\frac2{fk}\log\Gamma .
\]
The stored window $T_\omega^*\le T_\omega
\le(1+1/(fk))T_\omega^*$ says that
\[
 \delta_\omega:=\log T_\omega-\log T_\omega^*
          \in[0,\log(1+1/(fk))]\subseteq[0,1].
\]
In particular $|\delta_\omega-\delta_{\omega'}|\le1$.
The common $\Gamma$ term cancels before estimating, and GVC4 yields
\[
 |\log T_\omega-\log T_{\omega'}|
 \le\frac1{4k}+\frac23(1+\log k)+1
 \le4\log k.                                                 \tag{GVC.8}
\]
For the last inequality $k\ge5$ gives $\log k\ge1$ and
$1/(4k)\le1/4$; the resulting constant term
$1/4+2/3+1=23/12$ is at most $(10/3)\log k$.
All primes in (GVC.8) are those stored for the original two words.
One cannot obtain this bound from the separate size estimates
$\log T_\omega=O(m\ell)$; the fixed center factor and the
Vandermonde cancellation are essential.
% END GVC5

\subsection{The three clauses of the original count}

% BEGIN GVC6
\paragraph{GVC6: assemble the literal pointwise \texttt{Count\_at}.}
\label{claim:GVC6}
Assume the same eventual mass-ratio and base-volume inputs as
DSZ6, together with the original degree ties of DSZ3.
Set
\[
 C_Y=\max\{5,C_{YL}+1\},\qquad
 C_\epsilon=30+\frac{2+2C_Y}{\epsilon},\qquad
 C_c=\max\{C_\epsilon+2,\ 8+3C_Y,\ 4\}.                      \tag{GVC.9}
\]
For every sufficiently late member, define the original lower main
term
\[
 H=\sum_{j=1}^{m}\frac1f\sum_\tau
       \log\frac{w_j(D,\tau)v}{\pi}
       +m\log(g_0g_1).
\]
The QW6 lower estimate, applied with the exact DSZ6 mass and
factorial identity and DSZ4's checked coordinates, says
$X_\omega\ge H-C_\epsilon(\ell+m^2\ell/k)$.
DSZ9 proves $X_\omega-2\le\log T_\omega$ and
$\log T_\omega\le(8+3C_Y)m\ell$ on these same data.
We retain the full QW6 lower expression here, before the
coarser positive lower bound used for \texttt{TuSize}.

Since $\operatorname{rd}E\ge2$, $x>1$ and $\ell\ge1$,
\[
 \mathcal A=\log x+\ell+\log\operatorname{rd}E\ge\ell\ge1.
\]
Therefore
\[
 \begin{split}
 C_\epsilon(\ell+m^2\ell/k)+2
 &\le(C_\epsilon+2)(\mathcal A+m^2\mathcal A/k),\\
 H-C_c(\mathcal A+m^2\mathcal A/k)
 &\le\log T_\omega\le C_cm\ell .
 \end{split}\tag{GVC.10}
\]
Together with GVC5 and $C_c\ge4$ this is every clause of
$\mathrm{Count\_at}(\mathcal L,C_c)$.

The finite-index dictionary is exact. The source set
$\mathrm{SizeData}_2.\mathrm{lifts}(D)$ is $\{1,\ldots,m\}$;
QW6's zero-based $w_{j,\tau}$ is the original
$\mathrm{wj}(D,j+1,\tau)$. The two auxiliary primes are $g_0,g_1$,
and positivity gives
$\log g_0+\log g_1=\log(g_0g_1)=\mathrm{lgG}(D)$.
The real place set is the original $\mathrm{Tau}$ of cardinality
$f$, so its factor is exactly $1/f$, with no additional factor two.
These identifications turn (GVC.10) and (GVC.8) into the literal
source formula, rather than an estimate for a reindexed model.
% END GVC6

% BEGIN GVC7
\paragraph{GVC7: a complete uniform \texttt{CountF} on the same family.}
\label{claim:GVC7}
Suppose the original family is supplied with
\[
 \begin{split}
 &\mathrm{fam0}(\mathcal F,\mathrm{MassRatio}_2),\\
 &\mathrm{fam0}(\mathcal F,\mathrm{BaseVolumeEqualsMass}),\\
 &\mathrm{fam0}(\mathcal F,\mathrm{SigmaDegreeTwoSidedBoundOfGap}).
 \end{split}                                                 \tag{GVC.11}
\]
The third input is in fact available on every original family
by DSZ3, including its already recorded Harish--Chandra input.
The first two retain the same-record mass/volume hypotheses of
MF19 and RSL1--RSL3; no new realization is chosen here.
DSZ5--DSZ6 give the exact degree and mass identities after one
family threshold. DSZ4 and DSZ8--DSZ9 give all remaining
coordinate and prime-center bounds after finitely many further
thresholds. In particular, the row degree in the lower QW6
application is the one tied to this $D.\mathrm{dsig}$.

The number $C_{YL}$ is constant on this family: apply the
projection to the entry at index $19$ in the stored equality
$\mathrm{parConsts}(\mathcal F.\mathrm{params}(x,i))
 =\mathcal F.\mathrm{PK}$.
Thus choose the constants (GVC.9) first.
Then take a finite maximum of the mass/volume thresholds, all
DSZ/QW thresholds just listed, $2$, the threshold $m\ge44$,
and $\mathcal F.x_1+1$. For every $x$ beyond this maximum and
every $i\in\mathcal F.\mathrm{Ix}(x)$, the source equality
$\mathcal L.\mathrm{sc}.x=x$ is valid and GVC6 applies
simultaneously to all words. We have proved
\[
 \exists C_c\ge0\ \exists x_0\
 \forall x\ge x_0\ \forall i\in\mathcal F.\mathrm{Ix}(x),\
               \mathrm{Count\_at}(\mathcal F.L(x,i),C_c),
                                                               \tag{GVC.12}
\]
which is the original \texttt{CountF}.

This closes the count interface needed by the global ledger.
It does not assert the stronger
\texttt{FamilyVandermondeCountExplicit} or
\texttt{PartFiveFamilyBounds} package: their unused assertions are
not needed for (GVC.12). The proof is a direct human proof of the
literal consumer, with all three clauses checked.
% END GVC7

\subsection{Full numerical and weight conditions on the original family}

% BEGIN GVC8
\paragraph{GVC8: the complete original \texttt{Wts} predicate.}
\label{claim:GVC8}
Let the original $\mathrm{ParameterChoice}'(D,\mathcal L)$ be
given, $x\ge4$, $x\le X'$, and $C_n\ge3$.
For every word, the full source predicate holds:
\[
              \mathrm{Wts}(D,\omega,C_n,x,b,b',S,\eta).
\]

We check the clauses, including the appended weights that were
not needed for the shorter Hida-setting proof GLE1.
Put
\[
 \lambda_{\tau,j}=u_{\omega,\tau,j}
                  -\bigl((k-1)/2-j\bigr),\qquad
 z_{\tau,j}=w_j^{\rm raw}
              -\delta_j,\quad
 \delta_j=\frac{k}{2}+\frac{(2j+1)(t+1)}2,\quad0\le j<m,
                                                               \tag{GVC.13}
\]
where $w_j^{\rm raw}=\mathrm{wj}(D,j+1,\tau)$.
The original identity $D.\mathrm{hcW\_e}$ shows that this
$z_{\tau,j}$ is exactly \texttt{wJ}, with no word dependence
hidden in the raw appended coordinates. For $j<m$,
$0\le\delta_j\le N$ and
$\delta_{j+1}-\delta_j=t+1$.

The half-integrality of $u_{\omega,\tau,j}$ and evenness of $k$
make every $\lambda_{\tau,j}$ integral. Interlacing through the
integral coordinate $v_j$ gives
$u_{\omega,\tau,j}-u_{\omega,\tau,j+1}\ge1$:
writing the two half-integers as $r+1/2,s+1/2$, an integer
strictly between them forces $r\ge s+1$.
Hence $\lambda_{\tau,j}$ is nonincreasing.
GLE1, with its exact cloud and endpoint calculation, gives
\[
 \begin{array}{ll}
 \tau\text{ top}:&
 a\le\lambda_{\tau,0}\le2a,\quad
 b'\le-\lambda_{\tau,k-1}\le5b'/4<2b',\quad
 0\le\lambda_{\tau,j}\le S\ (0<j<k-1),\\
 \tau\text{ non-top}:&
 0\le\lambda_{\tau,j}\le3S\le C_nS\quad(0\le j<k).
 \end{array}                                                 \tag{GVC.14}
\]
That same parameter calculation proved
\[
 Nm^{12}\le w'',\quad S,\eta,mt\le w'',\quad
 \eta\ge1,\quad b'=xn w'',\quad
 N\le w''\le b'/x,\quad0<b'\le b.                            \tag{GVC.15}
\]
For clarity, the first inequality follows from the actual
$h_0$ bound $Nm^{11}\le\eta$ and
$m\eta\le w''$ supplied by the original \texttt{mteta}
inequality; the $S$ estimate follows from $kS\le w''$.
These are inequalities on the supplied parameters, not a new
choice of appended weights.

The two original Harish--Chandra gaps give
\[
 e_{j+1}-b_{j+1}\ge(t+3)/2,\qquad
 b_{j+1}-e_j\ge(t+3)/2.
\]
Thus $w_{j+1}^{\rm raw}-w_j^{\rm raw}\ge t+3$,
and subtracting $\delta_{j+1}-\delta_j=t+1$ proves
$z_{\tau,j+1}\ge z_{\tau,j}$.

At a top place the original window is
$[b_T+b'/x,\ b_T+b'/x+\eta]$.
Since $b_T\le5b'/4$ and both $b'/x,\eta\le b'/4$,
every $z_{\tau,j}\le7b'/4<2b'$.
At a thin place its raw window is $[w'',w''+\eta]$,
so $z_{\tau,j}\le2w''\le b'/2<2b'$.
These are all non-thick places, using the original disjoint
top/thick/thin partition.
At a thick place the raw window is $[b,b+\eta+1]$;
therefore
\[
                    b-N\le z_{\tau,j}\le b+\eta+1.
\]
For the first-weight clause, at top
$z_{\tau,0}\ge b_T+b'/x-N\ge b_T=-\lambda_{\tau,k-1}$.
Off top, the raw lower endpoint is at least $w''$,
so $z_{\tau,0}\ge w''-N\ge0\ge-\lambda_{\tau,k-1}$.

Finally all raw weights at a fixed place lie in one common
interval $[\mathrm{lo},\mathrm{lo}+\eta+1]$; choose
$\mathrm{lo}=b_T+b'/x$, $b$, or $w''$ in the three cases.
As $m\ge2$ and $t\ge1$,
\[
 z_{\tau,m-1}-z_{\tau,0}
 \le\eta+1-(m-1)(t+1)\le\eta.
\]
This proves every clause of the literal \texttt{Wts}, including
its last two inequalities, with the same $\omega,\tau$ and the
same original appended coordinates.
% END GVC8

% BEGIN GVC9
\paragraph{GVC9: the four large numerical inequalities.}
\label{claim:GVC9}
Let $N,m\ge1$, $1\le\ell\le m$, $c_L\ge3$,
$C_c\ge0$, $C_n\ge1$, and suppose
\[
 m\ge c_L+C_c+6C_n+9,\quad
 T_0\le C_cm\ell,\quad G_0\le2C_n\ell,\quad
 b'\ge Nm^4.
\]
Then
\[
 \begin{split}
 N^2(c_L\ell+T_0)&\le Nb',\\
 Nc_L\ell+T_0&\le b',\\
 m^2N^2(c_L\ell+T_0+G_0)&\le Nb'c_L\ell,\\
 (4C_n+9)m^3N&\le b'.
 \end{split}\tag{GVC.16}
\]
Indeed
$c_L\ell+C_cm\ell\le(c_L+C_c)m\ell
 \le m^2\ell\le m^3$.
The first inequality follows by multiplying by $N^2$ and using
$Nm^3\le Nm^4\le b'$.
For the second,
$Nc_L\ell+T_0\le N(c_L\ell+C_cm\ell)\le Nm^3\le b'$,
because $N\ge1$ and $C_cm\ell\ge0$.
Also $c_L+2C_n\le m$ and $C_c\le m$, so
\[
 c_L+C_cm+2C_n\le m+m^2\le2m^2\le c_Lm^2.
\]
Thus $c_L\ell+T_0+G_0\le c_Lm^2\ell$.
Multiplying by $m^2N^2$ and then using $Nm^4\le b'$
gives the third assertion.
Finally $4C_n+9\le m$, so the fourth follows from the same
$Nm^4$ bound. This proves the source's \texttt{nums\_alg}
with no suppressed asymptotic step.
% END GVC9

% BEGIN GVC10
\paragraph{GVC10: check all sixteen \texttt{Nums} clauses.}
\label{claim:GVC10}
For an original $\mathrm{ParameterChoice}'(D,\mathcal L)$
assume $x\ge4$, $x\le X'$, and
\[
 \begin{gathered}
 7\log x\le\ell,\quad
 n\ge3,\quad \log n\le3\log x,\quad
 \log q,\log d\le\ell,\quad
 \operatorname{rd}E\le2n^2\sqrt{qd},\\
 |\mathcal L.\mathrm{logDerivativeAtOne}|\le\ell,\quad
 \ell\le m,\quad m\ge c_L+C_c+6C_n+9,\\
 c_L\ge3,\quad C_c\ge0,\quad C_n\ge3,\quad
 C_c\le C_n,\quad C_{YL}+1\le C_n,\quad
 \log T_\omega\le C_cm\ell .
 \end{gathered}                                               \tag{GVC.17}
\]
Then the literal
$\mathrm{Nums}(D,T_\omega,g_0,g_1,C_n,x,b,c_L\ell,b',S,\eta)$
holds.

Here is a complete check in the original clause order.
First $b=e^\ell>0$.
Since $v=\operatorname{rd}E/2$ and the host discriminant bound
gives $v\ge1$,
\[
 \log v\le2\log n+\tfrac12(\log q+\log d)
          \le6\log x+\ell .
\]
It follows that
$\log(xbv)\le7\log x+2\ell\le3\ell\le c_L\ell$.
The third clause follows from
$\log T_\omega\le C_cm\ell\le C_nm(c_L\ell)$.
The original $g_Y$ window has $Y\le g_j\le2Y$ and
$\log Y\le C_{YL}\ell$, so
$\log g_j\le\log2+C_{YL}\ell\le(C_{YL}+1)\ell
 \le C_n(c_L\ell)$; this proves clauses four and five.

GVC8 gives $b'\le b=xb/x$, $\eta\ge1$, and
$S,\eta,mt\le w''$.
Also $S=4k\Delta\ge1$, since the original
$\Delta\ge Nm^{10}\ge1$ and $k\ge1$.
As $n\ge3$,
\[
              S+\eta+mt\le3w''\le n w''=b'/x .
\]
These prove clauses six through nine.
From $Nm^{12}\le w''\le b'$ and $m\ge1$ we have
$Nm^4\le b'$.
Set $T_0=\log T_\omega$ and
$G_0=\log(g_0g_1)\le2C_n\ell$ in GVC9.
Its first three bounds are exactly clauses ten through twelve.

The stored equality
$\mathcal L.\mathrm{logDerivative\_at\_one\_eq}$ identifies
the derivative in (GVC.17) with the original $\mathrm{dlogL}(D,1)$.
Hence
$|\mathrm{dlogL}(D,1)|\le\ell\le C_n(c_L\ell)$,
which is clause thirteen; this uses the supplied derivative
bound and does not prove a new analytic estimate.
Clauses fourteen and fifteen are $x\ge2$ and $v\ge1$.
The last GVC9 inequality is clause sixteen because
$1+2(2C_n+4)=4C_n+9$.
This completes the full conjunction, with the exact source
normalization $\widetilde L=c_L\ell$.
% END GVC10

% BEGIN GVC11
\paragraph{GVC11: one full \texttt{NumsL}/\texttt{WtsL} supplier, with constants before the threshold.}
\label{claim:GVC11}
Suppose the supplied family has an eventual uniform upper bound
\[
 \log T_\omega\le C_cm\ell
 \quad\text{and an eventual bound}\quad
 |\mathcal L.\mathrm{logDerivativeAtOne}|\le\ell .
                                                               \tag{GVC.18}
\]
The first is the second clause of \texttt{CountF};
\texttt{TuSize} also suffices for this particular input.
Write $\widehat C_c=\max(C_c,0)$ and fix
\[
 C_{n,0}=
 \max\bigl\{\max(3,\widehat C_c),\
                    \mathcal F.\mathrm{PK}[19]+1\bigr\}.
                                                               \tag{GVC.19}
\]
Then the exact source statement is
\[
 \boxed{\begin{gathered}
 \forall C_n\ge C_{n,0}\ \forall c_L\ge3\
 \exists x_0\ \forall x\ge x_0\ \forall i\in\mathcal F.\mathrm{Ix}(x):\\
 \mathrm{NumsL}(D,\mathcal L,C_n,c_L)
 \ \wedge\ \mathrm{WtsL}(D,\mathcal L,C_n).
 \end{gathered}}                                             \tag{GVC.20}
\]
Here \texttt{NumsL} is GVC10's \texttt{Nums} for every original
word, and \texttt{WtsL} is GVC8's \texttt{Wts} for every word.

We verify all uniformity and growth inputs.
For fixed $\epsilon>0$, $x^{14}\le e^{x^{\epsilon/3}}=X'$
eventually. One elementary proof puts $y=\log x$ and notes
$e^{(\epsilon/3)y}/y\longrightarrow\infty$, which follows from
$e^{(\epsilon/3)y}\ge(\epsilon y/3)^2/2$ for $y\ge0$.
This gives $14\log x\le x^{\epsilon/3}$ and hence
$7\log x\le\ell$ as well as $x\le X'$.

The actual \texttt{SizeM} exponent is
$1/2-\epsilon/16-\epsilon/100>2/5$; the source
\texttt{SizeN} exponents are $2+8\epsilon\pm\epsilon/100$.
Thus for $x\ge4$ the stored size inequalities imply
\[
 x^{2/5}\le m,\quad
 \ell\le x^{2/5},\quad n\ge x^2\ge3,\quad
 \log n\le3\log x.                                         \tag{GVC.21}
\]
The upper \texttt{SizeQD} exponents for $\log q$ and $\log d$
are $\epsilon/5+\epsilon/100<\epsilon$, so both are at most
$\ell$. The original discriminant formula gives
$\operatorname{rd}E\le2Q_S\sqrt{qd}\le2n^2\sqrt{qd}$;
its lower bound gives $\operatorname{rd}E\ge2$.
These are the already matched host bounds of AH16, on the same
host, together with $Q_S\le n^2$.

Now fix $C_n,c_L$ in the order in (GVC.20).
Take a threshold also making
\[
       x^{2/5}\ge c_L+\widehat C_c+6C_n+9.
\]
This is possible because $x^{2/5}$ tends to infinity.
The identity for the entry at index $19$ used in GVC7 gives
$C_{YL}+1\le C_n$ uniformly in $i$.
Enlarge the threshold to dominate the two bounds in (GVC.18),
the power-exponential threshold, $4$, and $\mathcal F.x_1+1$.
For every later $x,i$, GVC10 applies with $\widehat C_c$ and
GVC8 applies with $C_n$; both apply to all words at once.
This proves (GVC.20), including every clause of the literal
\texttt{nums\_wts\_of\_fam} and
\texttt{numsL\_of\_fam\_thr} interfaces.

In particular, choose $C_n=C_{n,0}$ and $c_L=3$ before the
last threshold. This gives the exact common numerical packet
used both by ordinary column slopes and by the adjacent level-prime
length calculation. This chapter does not replace $3$ by a
member-dependent constant.
% END GVC11

\subsection{Finite weight estimates and the original column slopes}

% BEGIN GVC12
\paragraph{GVC12: finite sums of the literal column weights.}
\label{claim:GVC12}
Let $D,\omega$ be arbitrary original data, and supply
\texttt{Nums} and \texttt{Wts} with parameters
$(C_n,x,b,\widetilde L,b',S,\eta)$.
Assume $x\ge K:=\max\{1,C_n\}$, and put $a=xb$.
For each original place define
\[
 s_\tau=\lambda_{\tau,k-1},\quad
 a_\tau=\sum_{j=0}^{k-1}\lambda_{\tau,j},\quad
 d_\tau=\sum_{j=0}^{k-1}(\lambda_{\tau,j}-s_\tau),\quad
 h_\tau=\lambda_{\tau,0}-s_\tau .
\]
These are respectively the last weight, \texttt{colAbs},
\texttt{colAbsP} and \texttt{firstColumnSpread}.
In particular $a_\tau$ is a signed sum, not a sum of absolute
values. By the integral-interlacing argument in GVC8 the weights
are nonincreasing, hence $d_\tau,h_\tau\ge0$.

The top place set has exactly $n_T$ members.
Explicitly, its bijection with $t\in D.\mathrm{host}.\mathrm{TT}$
is
$t\mapsto D.eTau^{-1}(\tau_1(t))$.
The original $\tau_1(t)$ lies in $\mathrm{Psi}$; injectivity
follows from the stored frame bijection, and membership in
\texttt{topPlaces}, defined as this image, gives surjectivity.
Thus this count is for the actual place set, not a doubled set
of complex embeddings.

The elementary consequences of \texttt{Nums} needed here are
\[
 b'\le b,\quad x\ge2,\quad v\ge1,\quad
 S,\eta\ge1,\quad S+\eta+mt\le b'/x,\quad
 N\widetilde L+\log T_\omega\le b'.                          \tag{GVC.22}
\]
The original ranks give $m\ge2,t\ge1,k\ge5,N\ge k$.
It follows that $b'\ge4x\ge8$, $a\ge16$, and
\[
 \widetilde L\ge\log(av)\ge\log a>1,\qquad
                  k\le b'\le a.
\]
Here $\log T_\omega\ge0$, since $T_\omega$ is the original
positive prime. Also $KS\le xS\le b'$.
Consequently the top spread is at most
$2a+2b'\le4a$, and every non-top spread is at most
$C_nS\le KS\le a$. Hence
\[
              \log(h_\tau+k)\le\log(5a)
                            \le2\widetilde L.              \tag{GVC.23}
\]
The last step uses $\log5\le\log a\le\widetilde L$.

At top, write $\lambda_{\tau,0}=a_T(\tau)$ and
$-\lambda_{\tau,k-1}=b_T(\tau)$.
The middle weights lie in $[0,S]$, so
\[
 a_\tau\le a_T(\tau)+kS,\qquad
 d_\tau\le2a+2kb'+kS,\qquad |s_\tau|\le2b'.
\]
Off top all entries lie in $[0,C_nS]$, giving
$a_\tau,d_\tau\le kKS$ and $|s_\tau|\le KS$.
Sum over the original $f$ places and use $K\ge1$:
\[
 \begin{split}
 \sum_\tau a_\tau
    &\le\sum_{\tau\ {\rm top}}a_T(\tau)+fkKS,\\
 \sum_\tau d_\tau
    &\le2n_Ta+fk(2b'+KS),\\
 \sum_\tau|s_\tau|&\le f(2b'+KS),\\
 \sum_\tau k^2\log(h_\tau+k)&\le2fk^2\widetilde L .
 \end{split}                                                 \tag{GVC.24}
\]
All summands multiplied by logarithms in the next step have
the signs explicitly checked here. This finite calculation
uses neither a square-mean estimate for arbitrary auxiliary
spaces nor the derivative clause of \texttt{Nums}.
% END GVC12

% BEGIN GVC13
\paragraph{GVC13: the exact \texttt{Conclusion5} and \texttt{Conclusion6} formulas.}
\label{claim:GVC13}
For the same $D$ suppose its original open lower and upper
column bounds hold, with arbitrary fixed upper coefficient $c_5$.
Set
\[
 K=\max\{1,C_n\},\qquad q_5=\max\{c_5,0\},\qquad
 C_-=K+2,\qquad C_+=(q_5+1)(K+4).                            \tag{GVC.25}
\]
For every $x\ge K$ these prove
$\mathrm{Conclusion5}(D,C_n,C_-,x)$ and
$\mathrm{Conclusion6}(D,C_n,C_+,x)$.

To prove this, take the arbitrary word, parameters, \texttt{Nums},
\texttt{Wts}, and original $\mathrm{CV}:\mathrm{ColVal}(D,\omega)$
quantified in those statements. Write
$M_\omega=(D.\mathrm{fm}\,\omega).\mathrm{lat}(0)$.
The lower raw inequality, proved for the actual model by GLS7, is
\[
 \mu(M_\omega)\ge
 -\frac{\log T_\omega}{k}\sum_\tau a_\tau
 -\log k\sum_\tau d_\tau
 -\sum_\tau k^2\log(h_\tau+k).
\]
Insert (GVC.24). Since $K+2$ dominates $K$, $2$ and the
coefficients of $2b'+KS$, this gives
\[
 \begin{split}
 \mu(M_\omega)\ge{}&
 -\sum_{\tau\ {\rm top}}\frac{a_T(\tau)}{k}\log T_\omega
 -C_-n_Ta\log k-C_-fk(b'+S)\log k\\
 &-C_-fS\log T_\omega-C_-fk^2\widetilde L .
 \end{split}                                                 \tag{GVC.26}
\]
This is the literal \texttt{Conclusion5}, including its signed
top sum and all four error terms.

The upper raw inequality is
\[
 \mu_{\max}(M_\omega)\le
 c_5\log(kv)\sum_\tau d_\tau
 +\sum_\tau k^2\log(h_\tau+k)
 +\log T_\omega\sum_\tau|s_\tau|.
\]
Here $\log(kv)\ge0$ because $k\ge5,v\ge1$.
We may therefore replace $c_5$ by $q_5$, and (GVC.24) yields
\[
 \begin{split}
 \mu_{\max}(M_\omega)\le{}&
 2q_5n_Ta\log(kv)
 +q_5fk(2b'+KS)\log(kv)\\
 &+f(2b'+KS)\log T_\omega
 +2fk^2\widetilde L .
 \end{split}
\]
The positive number $(q_5+1)(K+4)$ dominates
$2q_5$, $q_5(2+K)$, $2+K$ and $2$.
Thus
\[
 \begin{split}
 \mu_{\max}(M_\omega)\le{}&
 C_+n_Ta\log(kv)
 +C_+f(b'+S)\bigl(k\log(kv)+\log T_\omega\bigr)\\
 &+C_+fk^2\widetilde L ,
 \end{split}                                                 \tag{GVC.27}
\]
which is the literal \texttt{Conclusion6}.
For the actual GLS11 bound $c_5=4$, so $q_5=4$ and one may
take $C_+=5(K+4)$. The original coefficient $4$ is retained;
only the existential constant in the ensuing finite estimate
has been enlarged.

Equations (GVC.26)--(GVC.27) give independent human proofs of
the source's \texttt{ramifiedLattice\_\allowbreak of\_\allowbreak chain\_\allowbreak estimates}
and \texttt{classMap\_\allowbreak estimate\_\allowbreak of\_\allowbreak partial} interfaces.
They do not claim the stronger arbitrary-space
\texttt{zeta\_proposition\_second} or
\texttt{zeta\_proposition\_third}; the simpler direct sum
argument proves exactly the two original consumers required here.
% END GVC13

% BEGIN GVC14
\paragraph{GVC14: the ordinary column statements with the exact family quantifiers.}
\label{claim:GVC14}
Suppose
$\mathrm{FamilySlopeLowerOpenPart}(\mathcal F)$ and
$\mathrm{FamilySlopeMaxOpenPart}(\mathcal F,c_5)$ hold.
Fix any real $C_n$ first, and use the constants (GVC.25).
Take the maximum of the two eventual open-part thresholds and
$K=\max(1,C_n)$. For every later $x>\mathcal F.x_1$ and every
$i$, the original $\mathcal L.\mathrm{sc}.x=x$;
GVC13 therefore applies to every word, every tuple
$(b,\widetilde L,b',S,\eta)$ satisfying its numerical predicates,
and every supplied original \texttt{ColVal}.
This proves exactly
\[
 \begin{split}
 &\forall C_n\ \exists C_-,x_-\ \forall x\ge x_-\
   (x>\mathcal F.x_1)\ \forall i:
             \mathrm{Conclusion5}(D,C_n,C_-,x),\\
 &\forall C_n\ \exists C_+,x_+\ \forall x\ge x_+\
   (x>\mathcal F.x_1)\ \forall i:
             \mathrm{Conclusion6}(D,C_n,C_+,x).
 \end{split}                                                 \tag{GVC.28}
\]
These are the literal \texttt{ConclusionF5} and
\texttt{ConclusionF6}, not just their values at one convenient
tuple of numerical parameters.
GLS7 and GLS11 supply the two open inequalities on every original
$D$, conditional on the same \texttt{ColVal}; GLS12 supplies
their family formulations with $c_5=4$.
Their geometric and published inputs retain their stated scopes.
This step does not supply \texttt{ColVal}; the eventual
\texttt{tCV} input is used explicitly below.
% END GVC14

\subsection{Uniform ordinary slopes in the global ledger}

% BEGIN GVC15
\paragraph{GVC15: the two exact size dominations.}
\label{claim:GVC15}
For the original $\mathrm{SizesAt}(\mathcal L)$, suppose
$x\ge2$ and $x^{14}\le X'$. Then
\[
 \log k\le5\log x,\qquad
 fk(b'+S)\le n_Ta,\qquad fk^2\le n_Ta.                       \tag{GVC.29}
\]
Indeed \texttt{SizeN} and \texttt{SizeK}, with
$0<\epsilon<1/52$, give $n\le x^3$ and $k\le x^5$.
The positive lower \texttt{SizeTT} bound makes $n_T>0$;
as an integer it satisfies $n_T\ge1$.
The original \texttt{SizeDel} clause gives
$kS\le b'/(xn)$, and $x,n,k\ge1$, so $S\le b'$.
Using $f=2n$,
\[
 \begin{split}
 2fk&\le4x^8\le x^{14}\le X',\\
 fk(b'+S)&\le2fkb'\le X'b'=a\le n_Ta.
 \end{split}
\]
The middle comparison holds because $x^6\ge64\ge4$.
Also
\[
 fk^2\le2x^{13}\le x^{14}\le X'\le b\le a\le n_Ta.
\]
Here $X'\le b$ follows from $x^{\epsilon/3}\le x^\epsilon$.
Finally taking logarithms of $k\le x^5$ gives the first bound
of (GVC.29). This proves \texttt{jnt\_sizes\_slopes} with its
literal ranks and does not assume a new bound on $\Delta$.
% END GVC15

% BEGIN GVC16
\paragraph{GVC16: finite substitution into the exact \texttt{SlopesM\_at}.}
\label{claim:GVC16}
Suppose on this original $D,\mathcal L$ that
\texttt{Conclusion5} and \texttt{Conclusion6} hold with constants
$C_-,C_+$, that \texttt{NumsL}, \texttt{WtsL} hold with
$C_n,c_L$ where $c_L\ge0$, and that every word has an original
\texttt{ColVal}. Suppose also
$\log T_\omega\le C_cm\ell$ with $C_c\ge0$ and that (GVC.29)
holds. Put
\[
 P=\max\{C_-,C_+,0\},\qquad
                  C_s=2C_c+P(12+C_c+c_L).                  \tag{GVC.30}
\]
Then, for every word,
\[
 \boxed{\quad
 -C_sn_Ta\mathcal A(1+m/k)\le\mu(M_\omega),\qquad
 \mu_{\max}(M_\omega)\le C_sn_Ta\mathcal A .
 \quad}                                                      \tag{GVC.31}
\]
These are the two clauses of the original
$\mathrm{SlopesM\_at}(\mathcal L,C_s)$.

We check every numerical cost. Put $Z=n_Ta\mathcal A\ge0$.
The host bound $v\ge1$ implies $\operatorname{rd}E\ge2$, so
\[
 \ell\le\mathcal A,\quad
 0\le\log k\le5\mathcal A,\quad
 0\le\log(kv)\le6\mathcal A,\quad
 0\le\log T_\omega\le C_cm\mathcal A .
\]
The third inequality uses
$\log v\le\log\operatorname{rd}E$ and
$\mathcal A=\log x+\ell+\log\operatorname{rd}E$.
The top sum in (GVC.26) has exactly $n_T$ terms and
$a_T(\tau)\le2a$; its contribution is at most
$2C_c(m/k)Z$.

All four unsigned error terms of (GVC.26) are nonnegative.
Replacing their coefficient by $P$ is therefore valid even if
one of the original constants is negative.
The first two cost at most $5PZ$ each, by (GVC.29).
Since $fkS\le fk(b'+S)\le n_Ta$,
the third costs at most $PC_c(m/k)Z$.
The fourth, with $\widetilde L=c_L\ell$, costs at most $Pc_LZ$.
Thus the total lower cost is at most
\[
 [P(10+c_L)+(2C_c+PC_c)m/k]Z .
\]
The constant $C_s$ dominates both $P(10+c_L)$ and
$2C_c+PC_c$. This proves the lower assertion in (GVC.31).

For (GVC.27), its first term costs at most $6PZ$ and
the $k\log(kv)$ part of the second costs at most another $6PZ$.
The $\log T_\omega$ part costs at most $PC_c(m/k)Z\le PC_cZ$,
because $m\le k$. The last term costs at most $Pc_LZ$.
Their sum is bounded by $P(12+C_c+c_L)Z\le C_sZ$.
This proves the upper assertion.
In particular the exact source constant (GVC.30), including
the number $12$, is justified by a visible term-by-term calculation.
% END GVC16

% BEGIN GVC17
\paragraph{GVC17: the literal \texttt{SlopesMF} installation.}
\label{claim:GVC17}
Suppose this same original family has
\texttt{CountF}, \texttt{ConclusionF5}, \texttt{ConclusionF6},
the eventual derivative bound in (GVC.18), and
\[
 \exists x_{\rm CV}\ \forall x\ge x_{\rm CV}\
       \forall i\ \forall\omega,\
              \mathrm{Nonempty}(\mathrm{ColVal}(D,\omega)).
                                                               \tag{GVC.32}
\]
Extract a nonnegative $C_c$ and its threshold from
\texttt{CountF}. GVC11 supplies $C_n=C_{n,0}$ and $c_L=3$,
then a threshold for the full numerical and weight packet.
Only after fixing $C_n$ extract the two constants $C_-,C_+$
and their thresholds from (GVC.28).
Define the single $C_s\ge0$ by (GVC.30).

Take the maximum of the count, numerical, open-column,
$\mathrm{ColVal}$ and $x^{14}\le X'$ thresholds, of $2$, and
of $\mathcal F.x_1+1$. For each $x$ beyond it and each original
$i$, GVC15 gives the exact size inequalities.
For each word use an inhabitant of (GVC.32), on this same
$D,\omega$, in GVC16. Its constants and threshold do not depend
on the choice of inhabitant, the word, or the member.
We obtain
\[
 \exists C_s\ge0\ \exists x_0\
   \forall x\ge x_0\ \forall i\in\mathcal F.\mathrm{Ix}(x),\
          \mathrm{SlopesM\_at}(\mathcal F.L(x,i),C_s).
                                                               \tag{GVC.33}
\]
This is the original \texttt{SlopesMF}, and proves the full
\texttt{slopesMF\_of\_count\_fam} implication.
No finite or infinite metric is transported at this step:
$M_\omega$ has throughout been exactly
$(D.\mathrm{fm}\,\omega).\mathrm{lat}(0)$, with the original
norms in \texttt{ColVal}. Those norms and the compatible
column supplier are the retained GLS/RSL hypotheses.
% END GVC17

% BEGIN GVC18
\paragraph{GVC18: remove a redundant derivative supplier by a literal family projection.}
\label{claim:GVC18}
In the actual \texttt{FamP}, the separate eventual derivative
hypothesis of GVC11 and GVC17 is already supplied.
Indeed \texttt{Fm.sizes} gives
$\mathrm{SizesAt}(\mathcal F.L(x,i))$ for every
$x>\mathcal F.x_1$. The seventh conjunct of \texttt{SizesAt}
is \texttt{SizeQD}, whose sixth conjunct is precisely
\[
 |\mathcal L.\mathrm{logDerivativeAtOne}|
          \le\mathcal L.\mathrm{sc}.x^{
                         \mathcal L.\mathrm{sc}.\epsilon}.
\]
For $x\ge\mathcal F.x_1+1$ the original equalities
$\mathcal L.\mathrm{sc}.x=x$ and
$\mathcal L.\mathrm{sc}.\epsilon=\epsilon$ identify its right
side with $\ell=x^\epsilon$.
This proves the exact eventual hypothesis (GVC.18), with one
threshold before every $i$. It uses neither a new derivative
estimate nor a different host.

Consequently the exact outputs on a supplied original family are
\[
 \begin{gathered}
 \mathrm{CountF}\Longrightarrow
 \bigl[\text{the complete supplier (GVC.20)}\bigr],\\
 \mathrm{CountF}\wedge\mathrm{LowF}\wedge
 \mathrm{MaxF}(c_5)\wedge\mathrm{CVF}
       \Longrightarrow\mathrm{SlopesMF}.
 \end{gathered}                                               \tag{GVC.34}
\]
Here $\mathrm{LowF}$ and $\mathrm{MaxF}(c_5)$ abbreviate the
two original family open-part predicates, and $\mathrm{CVF}$
abbreviates (GVC.32).
For the selected original route, GLS12 supplies the open bounds
with $c_5=4$, and its \texttt{tCV} supplies (GVC.32) under the
same compatible \texttt{ColsTiedSym}/\texttt{adm3} data.
Combining those inputs with GVC7 therefore closes both count and
ordinary-slope entries of the one-family ledger.
The choice $c_L=3$ in the first line is also available to the
level-prime branch.

The word \emph{supplied} matters. This argument projects properties
from the existing \texttt{FamP}; it does not construct that family
or prove the analytic derivative estimate independently of its
size fields. The common selector, real-place second layer,
actual forms and mass/volume interfaces keep their inherited
conditional scope. The complementary quotient slopes, frame
estimate, loss sum and terminal \texttt{ChainTied} constructor
are separate remaining joins. No statement here closes the
whole RH or unrestricted Dirichlet-GRH human proof.
% END GVC18
